% GENERATED FILE. Edit canonical lessons, then run npm run generate:books.
% canonical-source-hash: fnv1a64:6d925d2299bafc97
% canonical-lessons: FR-C266-si, FR-C266-sil-pleut, FR-C266-ca-depend, FR-C266-revision-conditions

\chapter{If — si, s'il pleut, ça dépend}
\label{ch:fr-if}
\hlchaptermodality{266}

\begin{chapteropening}
\textbf{By the end of this chapter:} \emph{I can state a real condition with si and the present, and say that something depends, naming on what.}

Complete FR-C266-revision-conditions: a plan with a branch in it: a real condition with si, and an answer that depends.
\end{chapteropening}

\section[si, s'il]{si — if}

\label{lesson:FR-C266-si}



\begin{quote}
Say \textquotedblleft{}I don't know\textquotedblright{}. (\emph{Je ne sais pas.}) Then \textquotedblleft{}I have been living here for three years\textquotedblright{}. (\emph{J'habite ici depuis trois ans.})

Most sentences so far have said what is, or what was. This chapter says what follows \textbf{if} something else is true, and it begins with two letters.
\end{quote}

\begin{grammarlens}[title={Grammar Lens: si opens a condition}]
\textbf{si} — \textquotedblleft{}if\textquotedblright{}.

\begin{quote}
\textbf{Si} tu as faim, mange. — If you're hungry, eat. \textbf{Si} tu as froid, ferme la fenêtre. — If you're cold, shut the window.
\end{quote}

Whatever follows \emph{si} is the \textbf{condition}; the rest of the sentence is what follows from it. The condition can come first, with a comma, or second, without one: \emph{Ferme la fenêtre \textbf{si} tu as froid.} Notice that French \textbf{has} hunger and cold, \emph{tu as faim}, \emph{tu as froid}, where English \textbf{is} hungry and cold.
\end{grammarlens}

\begin{grammarlens}[title={Grammar Lens: s'il, and si for \textquotedblleft{}whether\textquotedblright{}}]
Before \textbf{il} and \textbf{ils}, \emph{si} drops its vowel: \textbf{s'il}, \textbf{s'ils}. Before \emph{elle} it keeps it: \emph{si elle}. You have said \emph{s'il} since the chapter on please, because \emph{s'il vous plaît} is \textquotedblleft{}if it pleases you\textquotedblright{}.

After a verb of knowing or asking, \emph{si} also means \textquotedblleft{}whether\textquotedblright{}, and your \emph{je ne sais pas} comes in useful:

\begin{quote}
\emph{\textbf{Je ne sais pas si} c'est loin.} — I don't know whether it's far.
\end{quote}
\end{grammarlens}

\begin{cousinweb}
You have met \textbf{si} before, as the \textquotedblleft{}yes\textquotedblright{} that contradicts a negative: \emph{Tu ne viens pas ? — Si !} That one goes back to Latin \emph{sīc}, \textquotedblleft{}so, thus\textquotedblright{}. The \emph{si} of this lesson is inherited from Latin \emph{sī}, \textquotedblleft{}if\textquotedblright{}, the Latin word that opened a condition (Lewis \& Short, under \emph{si}). EtymDB 2.1 lists both Latin words behind the French spelling. Two words, worn into one shape.

The ear tells them apart by their place. The \textquotedblleft{}yes\textquotedblright{} \emph{si} stands alone, as an answer. The \textquotedblleft{}if\textquotedblright{} \emph{si} opens a whole clause.
\end{cousinweb}

\subsection*{Your turn}
Say these aloud:

\begin{itemize}
\raggedright
  \item \emph{Si tu as faim, mange.}
  \item the same, condition second — \emph{Mange, si tu as faim.}
  \item \emph{Si tu as froid, ferme la fenêtre.}
  \item \emph{Je ne sais pas si c'est loin.}
\end{itemize}

\begin{tcolorbox}[breakable,colback=teal!4,colframe=teal!35!black,title={Before you move on}]
Which half of \emph{si tu as faim, mange} is the condition? (\emph{\textbf{Si tu as faim.}}) What does \emph{si} become before \emph{il}? (\emph{\textbf{S'il.}}) What does \emph{si} mean after \emph{je ne sais pas}? (\textbf{Whether}.) Which Latin words are behind the two kinds of \emph{si}? (\emph{\textbf{Sī}}, \textquotedblleft{}if\textquotedblright{}, and \emph{\textbf{sīc}}, \textquotedblleft{}so\textquotedblright{}.)
\end{tcolorbox}
\section[s'il pleut, je reste]{s'il pleut, je reste — the condition that really happens}

\label{lesson:FR-C266-sil-pleut}



\begin{quote}
Say \textquotedblleft{}if\textquotedblright{}. (\emph{Si.}) Say \textquotedblleft{}it was raining\textquotedblright{}. (\emph{Il pleuvait.}) Now the present: \textquotedblleft{}it's raining\textquotedblright{}. (\emph{Il pleut.})
\end{quote}

\begin{grammarlens}[title={Grammar Lens: both halves in the present}]
\begin{quote}
\textbf{S'il pleut, je reste} à la maison. — If it rains, I stay at home. \textbf{Si le train est cher, je prends} la voiture. — If the train is expensive, I take the car.
\end{quote}

This is a \textbf{real} condition, true now or whenever it happens. Both halves take the plain present, as in English: nothing new to build.

The result half can also be advice, with the \emph{tu devrais} of chapter two hundred and sixty-one, or a command:

\begin{quote}
\emph{Si tu es malade, \textbf{tu devrais} dormir.} — If you're ill, you should sleep. \emph{Si tu es fatigué, \textbf{dors} !} — If you're tired, sleep!
\end{quote}
\end{grammarlens}

\begin{culture}
A real condition is a plan with a branch in it. Someone who takes the train to work every day lives by two of them:

\begin{quote}
\emph{Si je rate le train, je prends le métro. Si je rate le métro, je reste à la maison.}
\end{quote}

Compare \emph{donc}. \emph{Il pleut, donc je reste} says it \textbf{is} raining. \emph{S'il pleut, je reste} says it \textbf{might}, and here is what happens if so. Same consequence; the second has not decided whether the cause is true.
\end{culture}

\begin{grammarlens}[title={Grammar Lens: tomorrow, still in the present}]
For one particular occasion still to come, the result can move into the future, with the \emph{aller} + infinitive you know. The \emph{si} half does not move:

\begin{quote}
\emph{S'il pleut \textbf{demain}, je \textbf{vais rester} à la maison.} — If it rains tomorrow, I'm going to stay at home.
\end{quote}

\emph{Il pleut}, the present, even about tomorrow. English does the same, \textquotedblleft{}if it rains tomorrow\textquotedblright{}, so here the ear is on your side.
\end{grammarlens}

\begin{cousinweb}
\textbf{rester} goes back to Latin \emph{restāre} (EtymDB 2.1), \textquotedblleft{}to stay behind, to hold still\textquotedblright{} (Lewis \& Short, under \emph{resto}): \emph{re-}, \textquotedblleft{}back\textquotedblright{}, plus \emph{stāre}, \textquotedblleft{}to stand\textquotedblright{}, the standing verb behind \emph{j'étais}. To stay is to stand still while the others go on.
\end{cousinweb}

\subsection*{Your turn}
Say these aloud:

\begin{itemize}
\raggedright
  \item \emph{S'il pleut, je reste à la maison.}
  \item the branch — \emph{Si je rate le train, je prends le métro.}
  \item advice in the result — \emph{Si tu es malade, tu devrais dormir.}
  \item tomorrow — \emph{S'il pleut demain, je vais rester à la maison.}
\end{itemize}

\begin{tcolorbox}[breakable,colback=teal!4,colframe=teal!35!black,title={Before you move on}]
Which tense goes in the \emph{si} half of a real condition? (The \textbf{present}, even for tomorrow.) What is the difference between \emph{il pleut, donc je reste} and \emph{s'il pleut, je reste}? (The first says it \textbf{is} raining; the second has \textbf{not decided}.) What is \emph{rester}, taken apart? (\emph{\textbf{Re-}} and \emph{\textbf{stāre}}: to stand back.)
\end{tcolorbox}
\section[ça dépend]{ça dépend — the condition you do not have to finish}

\label{lesson:FR-C266-ca-depend}



\begin{quote}
Say \textquotedblleft{}if it rains, I stay at home\textquotedblright{}. (\emph{S'il pleut, je reste à la maison.})

That sentence names its condition. Sometimes you want the condition without saying what it is.
\end{quote}

\subsection*{You'll want to know: ça dépend}
\textbf{ça dépend} — \textquotedblleft{}it depends\textquotedblright{}.

\begin{quote}
— \emph{Tu vas au cinéma demain ?} — \emph{\textbf{Ça dépend.}} — Going to the cinema tomorrow? It depends.
\end{quote}

You have \textbf{dépendre}, \textquotedblleft{}to depend\textquotedblright{}. \emph{Ça dépend} is its \textquotedblleft{}it\textquotedblright{} form, and on its own it is a complete answer: it says a condition exists and declines to say which.

\emph{Ça} is the everyday \textquotedblleft{}it, that\textquotedblright{}, and \emph{\textbf{tout dépend}} is \textquotedblleft{}it all depends\textquotedblright{}.

\begin{grammarlens}[title={Grammar Lens: ça dépend de}]
Add \textbf{de} and you can name the condition after all:

\begin{quote}
\emph{\textbf{Ça dépend du} temps.} — It depends on the weather. \emph{\textbf{Ça dépend de} toi.} — It depends on you. \emph{\textbf{Ça dépend de l'}heure.} — It depends on the time.
\end{quote}

\emph{De} and \emph{le} weld into \textbf{du}, as in \emph{content du voyage}. English says \emph{depend \textbf{on}}; French says \emph{dépend \textbf{de}}, \textquotedblleft{}hangs \textbf{from}\textquotedblright{}, which is the older picture.

Or name it with a whole \emph{si}: \emph{Ça dépend. S'il pleut, je reste à la maison.} The \emph{ça dépend} announces a condition; the \emph{si} sentence states it.
\end{grammarlens}

\begin{cousinweb}
\textbf{dépendre} has no parent in EtymDB 2.1; it is usually traced to Latin \emph{dēpendēre}, \textquotedblleft{}to hang down from something\textquotedblright{} (Lewis \& Short, under \emph{dependeo}): \emph{dē}, \textquotedblleft{}down from\textquotedblright{}, and \emph{pendēre}, \textquotedblleft{}to hang\textquotedblright{}. To depend on something is to \textbf{hang from} it, and the answer swings until the condition settles. The same \emph{pendēre} is inside \emph{pendant}, which you met with \emph{pendant que}. Its close relative \emph{pendere}, \textquotedblleft{}to make hang, to weigh on the scales\textquotedblright{} (Lewis \& Short, under \emph{pendo}), is the one behind \emph{penser}. English \textbf{pendant} and \textbf{pending} belong to the family too.
\end{cousinweb}

\subsection*{Your turn}
Say these aloud:

\begin{itemize}
\raggedright
  \item answer \emph{Tu vas au cinéma demain ?} — \emph{Ça dépend.}
  \item name it — \emph{Ça dépend du temps.}
  \item name it with a whole condition — \emph{Ça dépend. S'il pleut, je reste à la maison.}
  \item hand it back — \emph{Ça dépend de toi. Qu'est-ce que tu en penses ?}
\end{itemize}

\begin{tcolorbox}[breakable,colback=teal!4,colframe=teal!35!black,title={Before you move on}]
What does \emph{ça dépend} say, and what does it not say? (That there \textbf{is} a condition — not \textbf{which}.) What do you add to name it? (\emph{\textbf{De}}: \emph{ça dépend de toi}.) What is hanging, in \emph{dēpendēre}? (\textbf{The answer}, from the condition.)
\end{tcolorbox}
\section[(review)]{Review — if, and it depends}

\label{lesson:FR-C266-revision-conditions}



\begin{quote}
Say \textquotedblleft{}if he is tired\textquotedblright{}, with the vowel \emph{si} loses. (\emph{S'il est fatigué.}) Then \textquotedblleft{}it depends\textquotedblright{}. (\emph{Ça dépend.})
\end{quote}

\begin{grammarlens}[title={What you've built}]
Three lessons, and a way of talking about what has not happened yet. The condition, opened by \emph{si}: \emph{s'il pleut}. The result, in the present, the future with \emph{aller}, or a piece of advice. And the condition not yet named: \emph{ça dépend}, or named after all: \emph{ça dépend du temps}.

\begin{quote}
— \emph{Tu vas à la plage demain ?} — \emph{Ça dépend du temps. S'il fait beau, je vais à la plage. S'il pleut, je vais rester à la maison.}
\end{quote}

Two branches, one for each kind of weather, and the \emph{si} half in the present both times, although both are about tomorrow.
\end{grammarlens}

\subsection*{Your turn}
Say these aloud:

\begin{itemize}
\raggedright
  \item \emph{Ça dépend du temps.}
  \item the first branch — \emph{S'il fait beau, je vais à la plage.}
  \item the second — \emph{S'il pleut, je vais rester à la maison.}
  \item advice in a branch — \emph{Si tu es fatigué, tu devrais dormir.}
  \item your own plan for tomorrow, with \emph{ça dépend} and two branches
\end{itemize}

\begin{tcolorbox}[breakable,colback=teal!4,colframe=teal!35!black,title={Before you move on}]
In \emph{s'il fait beau demain}, which tense is \emph{fait}? (The \textbf{present}.) What does \emph{si} become before \emph{il}? (\emph{\textbf{S'il.}}) And what does \emph{ça dépend} leave unsaid? (\textbf{Which} condition.)
\end{tcolorbox}
